\chapter{Limites, risques et perspectives}

\section{Limites actuelles}
Malgre les progres architecturaux, plusieurs limites subsistent :
\begin{itemize}
\item couverture encore partielle des familles d'indicateurs ;
\item sensibilite potentielle a la qualite des donnees sur certains titres peu liquides ;
\item profondeur encore limitee des tests de robustesse multi-regimes.
\end{itemize}

\section{Risques methodologiques}
Le Tableau~\ref{tab:risk-matrix} propose une matrice synthetique des risques.

\begin{table}[H]
\centering
\caption{Matrice des risques methodologiques}
\label{tab:risk-matrix}
\begin{tabular}{p{4.2cm}p{3cm}p{6.3cm}}
\toprule
\textbf{Risque} & \textbf{Niveau} & \textbf{Mitigation prioritaire} \\
\midrule
Overfitting residuel & Moyen/Eleve & Renforcer protocoles holdout et tests de permutation. \\
Drift de regime & Moyen & Ajouter un module de detection de regime plus strict. \\
Biais de donnees & Moyen & Consolider controles de completude/fraicheur par symbole. \\
Complexite parametrique & Moyen & Limiter l'espace de recherche et auditer les degres de liberte. \\
\bottomrule
\end{tabular}
\end{table}

Le Tableau~\ref{tab:risk-matrix} montre que les risques majeurs sont
methodologiques avant d'etre purement logiciels.

\section{Plan d'amelioration priorise}
Le plan d'evolution est organise en trois horizons :
\begin{enumerate}
\item \textbf{Court terme} : renforcer la gouvernance des experiences
(protocoles de benchmark, versions de datasets, checks automatiques).
\item \textbf{Moyen terme} : etendre les familles d'indicateurs en conservant le
meme cadre OOS et les memes contraintes de robustesse.
\item \textbf{Long terme} : connecter la plateforme a des flux de donnees
institutionnels et formaliser une chaine pre-production.
\end{enumerate}

\section{Perspective scientifique}
La perspective principale est de maintenir une architecture qui se complexifie
sans perdre ses invariants. Toute extension doit conserver :
\begin{itemize}
\item la separation des responsabilites,
\item la reproductibilite des verdicts,
\item la priorite donnee a la robustesse OOS.
\end{itemize}

